\chapter{SOLUCIÓN PLANTEADA}
\label{cap:solucion}

\section{FORMULACIÓN Y DESCRIPCIÓN DE LA SOLUCIÓN}
\label{sec:solucion}

El problema de ruteo de censos se relaciona con dos variantes conocidas de la investigación
operativa. El \textit{Multiple Travelling Salesman Problem} (mTSP) generaliza el TSP a múltiples
agentes \citep{bektas2006multiple}. El \textit{Vehicle Routing Problem with Time Windows} (VRPTW)
introduce restricciones de ventanas horarias en nodos \citep{solomon1987algorithms,
toth2014vehicle}. Su resolución práctica recurre a metaheurísticas como \textit{Guided Local Search} (GLS)
\citep{voudouris1999guided}.

Las herramientas de ruteo de propósito general resuelven variantes del problema de ruteo de
vehículos; sin embargo, ninguna cubre la integración con el flujo de trabajo de este proyecto:
partir del catastro que dejaron
las iteraciones anteriores de Arbocensus, publicar el plan como asignación de trabajo a censistas
identificados y recoger la observación histórica de cada árbol. Esa
cadena ---catastro legado, optimización, asignación, captura en terreno e historial--- es lo que el
producto debe aportar, y es específica del dominio y no del algoritmo.

Se propone el desarrollo de una plataforma web que formule la planificación del censo como un
ruteo de vehículos de rutas abiertas con visitas opcionales penalizadas y restricción de
duración por ruta, sobre distancias caminables de la red vial real y no euclidianas. La plataforma
debe resolver esa formulación mediante programación con restricciones y cubrir las cuatro
condiciones del ruteo urbano real que el mTSP y el VRPTW no reúnen juntas:
\begin{enumerate}
  \item determinar automáticamente el número de rutas necesarias $k$ dado el conjunto de árboles $P =
        \{p_1, \ldots, p_n\}$ y los parámetros $T_{\min}$, $T_{\max}$, $s$ (tiempo de servicio por
        árbol), sin intervención manual;
  \item respetar restricciones de jornada global por ruta, no solo ventanas en nodos;
  \item generar rutas abiertas, adecuadas para operaciones en terreno;
  \item atender simultáneamente el número de rutas, el tiempo total de desplazamiento y el
        desbalance de carga entre equipos.
\end{enumerate}

Se consideró también un enfoque \textit{route-first, cluster-second} \citep{beasley1983route}:
resolver un único TSP sobre
los $n$ árboles y cortar el recorrido resultante en tramos de duración acumulada $\leq T_{\max}$. Se
descartó porque el corte por duración no controla el balance ni la compacidad de cada tramo, y
porque los extremos de cada tramo quedan fijados por un recorrido único que ignora, al construirse,
la partición que luego se le impone.

El sistema no se agota en el motor de optimización: debe cubrir el ciclo operacional completo de
una campaña, desde la constitución del conjunto de árboles a censar hasta el seguimiento del avance
en terreno. El administrador importa el arbolado, configura los parámetros de jornada,
ejecuta la optimización, publica el plan resultante y asigna cada ruta a un censista. El censista
recorre su ruta desde el navegador de su teléfono y registra, parada por parada, el estado en que
encontró cada árbol. Cada registro es una observación histórica.

\section{USUARIOS OBJETIVO}
\label{sec:usuarios}

La plataforma se dirige a dos perfiles de usuario con contextos de uso distintos, diferencia que
gobierna las decisiones de diseño del Capítulo~\ref{cap:diseno}.

\begin{enumerate}
  \item Administrador de campaña. Coordina el censo dentro del proyecto Arbocensus. No se supone en
        él conocimiento de optimización combinatoria: no debe declarar cuántos equipos
        necesita ni ajustar coeficientes del \textit{solver}. Trabaja desde un
        computador de escritorio, con conexión estable, en sesiones de planificación previas a la
        jornada de terreno. Se espera un único administrador activo por campaña.
  \item Censista. Participa en la jornada de terreno recorriendo a pie la ruta que le fue asignada.
        Utiliza su propio teléfono con datos móviles, en la vía pública, durante una jornada acotada
        a dos o tres horas, extensión fijada como definición del proyecto por la autonomía de la
        batería del dispositivo. La conectividad puede degradarse de forma intermitente.
\end{enumerate}

De ese contraste se derivan tres exigencias concretas: el portal debe ser operable con una sola mano
en la pantalla de un teléfono, sin depender de un formato de escritorio; debe recorrerse de árbol en
árbol sin que el censista desbloquee el dispositivo en cada parada; y el registro de una visita debe
sobrevivir a una pérdida momentánea de red en lugar de perderse.

\section{REQUISITOS FUNCIONALES}
\label{sec:requisitos-funcionales}

La Tabla~\ref{tab:requisitos} enumera los requisitos que la solución compromete, agrupados por la
etapa de la cadena operacional de la Sección~\ref{sec:solucion} a la que pertenecen.

\begin{xltabular}{\textwidth}{@{}l X@{}}
  \caption{Requisitos funcionales de la solución.}
  \label{tab:requisitos} \\
  \toprule
  ID & Requisito \\
  \midrule
  \endfirsthead
  \toprule
  ID & Requisito \\
  \midrule
  \endhead
  \bottomrule
  \endlastfoot
  \multicolumn{2}{@{}l}{\textit{Integración del catastro existente}} \\
  \addlinespace
  RF-01 & Consultar el arbolado registrado en las bases de datos legadas del proyecto, con sus áreas
          de censo. \\
          \addlinespace
  RF-02 & Delimitar territorialmente el arbolado que se censará en una campaña. \\
  \addlinespace
  RF-03 & Constituir el conjunto de árboles de una campaña (\textit{dataset}) a partir de esa
          delimitación, junto con el historial de observaciones que cada árbol ya tenga. \\
          \addlinespace
  RF-25 & Filtrar la exploración legada por estado y por \textit{dataset} de origen, y reincorporar
          un árbol ya presente en otro \textit{dataset} cuando la campaña lo requiere. \\
          \addlinespace
  \multicolumn{2}{@{}l}{\textit{Planificación del censo}} \\
  \addlinespace
  RF-04 & Configurar los parámetros de la jornada: duración mínima $T_{\min}$, duración máxima
          $T_{\max}$ y tiempo de servicio por árbol. \\
          \addlinespace
  RF-05 & Estimar cuántos censistas exigiría una campaña antes de planificarla. \\
  \addlinespace
  RF-06 & Asignar cada árbol del \textit{dataset} a una ruta y fijar el orden de visita, sobre los
          tiempos de desplazamiento a pie que impone la red vial real. \\
          \addlinespace
  RF-07 & Determinar el número de rutas $k$ como salida de la planificación, sin que el
          administrador lo declare. \\
          \addlinespace
  RF-08 & Reportar los árboles que la planificación no logró incorporar a ninguna ruta. \\
  \addlinespace
  RF-09 & Exponer las métricas de un plan ---cobertura, número de rutas, tiempo de desplazamiento y
          balance de carga--- antes de que se publique. \\
          \addlinespace
  RF-10 & Presentar las rutas de un plan sobre el trazado que el censista recorrerá a pie. \\
  \addlinespace
  RF-22 & Comprobar la factibilidad de la configuración antes de lanzar la optimización,
          distinguiendo bloqueos de advertencias. \\
          \addlinespace
  RF-23 & Recomendar automáticamente la solución preferible entre las disponibles de un
          \textit{dataset}. \\
          \addlinespace
  \multicolumn{2}{@{}l}{\textit{Publicación y asignación}} \\
  \addlinespace
  RF-11 & Publicar un plan como vigente, con a lo más un plan vigente por \textit{dataset}. \\
  \addlinespace
  RF-12 & Asignar y desasignar censistas a las rutas del plan vigente. \\
  \addlinespace
  RF-24 & Reportar la carga acumulada de cada censista y proponer una asignación de rutas que la
          equilibre en el largo plazo. \\
          \addlinespace
  \multicolumn{2}{@{}l}{\textit{Captura en terreno}} \\
  \addlinespace
  RF-13 & Entregar a cada censista las rutas que le fueron asignadas, y solo esas. \\
  \addlinespace
  RF-14 & Presentar las paradas de una ruta en orden secuencial estricto. \\
  \addlinespace
  RF-15 & Registrar en cada parada la observación del árbol ---su estado, su fotografía y sus
          notas---, fechada y asociada al ejemplar. \\
          \addlinespace
  RF-16 & Registrar la omisión de una parada con un motivo obligatorio. \\
  \addlinespace
  RF-17 & Conservar el historial completo de observaciones de cada árbol. \\
  \addlinespace
  RF-18 & Preservar el registro de una visita cuando el dispositivo pierde la conexión. \\
  \addlinespace
  \multicolumn{2}{@{}l}{\textit{Seguimiento y administración}} \\
  \addlinespace
  RF-19 & Presentar el avance del censo: árboles censados, omitidos y pendientes. \\
  \addlinespace
  RF-20 & Desglosar el avance por ruta y por censista. \\
  \addlinespace
  RF-21 & Administrar las cuentas de usuario y sus roles, sin que dar de baja a una persona borre el
          trabajo que registró. \\
          \addlinespace
  RF-26 & Editar los metadatos de un \textit{dataset} (nombre y descripción) y desactivar árboles
          sin borrarlos del sistema. \\
\end{xltabular}

\section{ATRIBUTOS DE CALIDAD}
\label{sec:atributos-calidad}

Los atributos de calidad comprometidos son los cuatro presentados en laTabla~\ref{tab:atributos-calidad}.
\begin{table}[H]
  \centering
  \caption{Atributos de calidad comprometidos por la solución.}
  \label{tab:atributos-calidad}
  \small
  \begin{tabularx}{\textwidth}{@{}p{3.1cm} l X@{}}
    \toprule
    Atributo &
    Prioridad &
    Exigencia sobre el sistema \\
    \midrule
    {\raggedright Robustez del plan\par} &
    Alta &
    El plan sigue siendo utilizable cuando cambian las condiciones de la campaña ---el tiempo de
    servicio por árbol, la duración de la jornada y el conjunto de árboles---, entendiendo por
    utilizable que ningún árbol queda fuera y que ningún censista carga con un trabajo
    desproporcionado respecto de sus pares. \\
    {\raggedright Operabilidad en terreno\par} &
    Alta &
    El portal del censista se opera a pie, en la vía pública, con un teléfono como único dispositivo
    y una conectividad intermitente, sin que esas condiciones comprometan lo que se registra. \\
    {\raggedright Eficiencia de cómputo\par} &
    Media &
    Planificar una campaña completa cuesta un tiempo compatible con la preparación de una jornada de
    terreno, y la calidad del plan no depende de conceder al cómputo un presupuesto mayor. \\
    {\raggedright Calidad geográfica de las rutas\par} &
    Media &
    Cada ruta es recorrible a pie sin retrocesos ni cruces sobre sí misma, juzgada sobre el trazado
    real de calles y no sobre distancias en línea recta. \\
    \bottomrule
  \end{tabularx}
\end{table}

La robustez del plan es alta porque el sistema no se opera bajo una única configuración: los dos
parámetros que el administrador fija en cada campaña desplazan el problema entre regímenes de
trabajo distintos, de modo que un plan correcto bajo una combinación no acredita nada bajo otra.
Sus dos componentes son los que un censo no puede sacrificar: cobertura y reparto. El reparto se
observa mediante el índice de balance de carga ---el cociente entre la duración de la ruta más
corta y la de la más larga---, la única noción de balance que emplea esta memoria; el umbral
exigible se fija en CA-03 como exigencia de reparto, no como resultado de una medición previa.

La operabilidad en terreno eleva a propiedad del producto las tres exigencias del contexto del
censista; ninguna es una función que el sistema ofrezca, sino una condición bajo la que todas sus
funciones deben seguir sirviendo. Su prioridad es alta porque una observación perdida en terreno no
se recupera: la jornada ya terminó y el árbol quedó sin dato.

La eficiencia de cómputo aparece con prioridad media porque la planificación no es una interacción
sino un encargo: la magnitud relevante es el costo de planificar una campaña, no la latencia
percibida. Se añade una exigencia de suficiencia: como la búsqueda se corta por tiempo y no por
convergencia demostrada, el presupuesto asignado debe capturar la mejora alcanzable, lo que se
observa comparando contra un presupuesto mayor.

La calidad geográfica de las rutas se juzga sobre el trazado real de calles porque el censista
camina por la vereda y no en línea recta. Su prioridad es media porque una ruta con auto-cruces
sigue siendo ejecutable ---no invalida el censo, lo encarece--- y porque el tiempo de
desplazamiento, objetivo explícito del modelo, ya penaliza buena parte de esa
ineficiencia.\footnote{Los cuatro atributos se retoman en la
Sección~\ref{sec:cumplimiento-atributos}.}

\section{RESTRICCIONES}
\label{sec:restricciones}

La Tabla~\ref{tab:restricciones} enumera las restricciones bajo las que se desarrolla el producto. 

\small
\begin{xltabular}{\textwidth}{@{}l X@{}}
  \caption{Restricciones consideradas en la elaboración del producto.}
  \label{tab:restricciones} \\
  \toprule
  ID & Restricción \\
  \midrule
  \endfirsthead
  \toprule
  ID & Restricción \\
  \midrule
  \endhead
  \bottomrule
  \endlastfoot
    
    \addlinespace
    RE-03 &
    Las bases de datos legadas pertenecen a sistemas ajenos al alcance de esta iteración y se
    acceden estrictamente en modo de solo lectura. \\
    \addlinespace
    RE-06 &
    La conectividad en terreno es intermitente, de modo que el portal del censista no puede suponer
    red disponible en el instante del registro. \\
\end{xltabular}
\normalsize


\section{CRITERIOS DE ACEPTACIÓN}
\label{sec:criterios-aceptacion}

Los criterios de aceptación son la traducción de los atributos de calidad a condiciones
verificables.

\begin{table}[H]
  \centering
  \begin{threeparttable}[b]
  \caption{Criterios de aceptación comprometidos para la solución.}
  \label{tab:criterios}
  \small
  \begin{tabularx}{\textwidth}{@{}l X X@{}}
    \toprule
    ID &
    Criterio &
    Valor exigido \\
    \midrule
    CA-01 &
    Cobertura del censo. &
    Ningún árbol del \textit{dataset} queda fuera del plan. \\
    \addlinespace
    CA-02 &
    Respeto de la jornada máxima. &
    Ninguna ruta excede $T_{\max}$. \\
    \addlinespace
    CA-03 &
    Balance de carga entre equipos.$^{1}$ &
    Índice de al menos $0{,}80$. \\
    \addlinespace
    CA-04 &
    Determinación de la flota. &
    El número de rutas $k$ es una salida del \textit{solver} y no un dato que el administrador
    declare. \\
    \addlinespace
    CA-05 &
    Costo de planificar una campaña.$^{2}$ &
    Menos de 5 minutos. \\
    \addlinespace
    CA-06 &
    Mejora frente a la alternativa trivial. &
    Reducción estrictamente positiva del tiempo total de desplazamiento frente a un heurístico de
    vecino más cercano, sobre la misma instancia y la misma matriz de costos. \\
    \bottomrule
  \end{tabularx}
  \begin{tablenotes}[flushleft]\footnotesize
    \item[1] Medido como el cociente entre la duración de la ruta más corta y la de la más larga.
    \item[2] Tiempo de reloj de una corrida completa del proceso de optimización, con la matriz de
      costos en caché. Comprende el presupuesto de tiempo que se configura al \textit{solver}, de
      modo que el criterio acota esa configuración además de su resultado.
  \end{tablenotes}
  \end{threeparttable}
\end{table}
